%! Semi-log Plots — Unit 2, Lesson 2.15 — Student Packet Cover Sheet
\documentclass[10pt]{article}
\usepackage{precalculus-article}
\usepackage{precalculus-boxes}
\usepackage{ltablex}
\keepXColumns

\begin{document}

% Full-bleed navy banner
\begin{tikzpicture}[remember picture, overlay]
  \fill[navy] (current page.north west)
    rectangle ([yshift=-0.9in]current page.north east);
\end{tikzpicture}
\vspace{-0.8in}

\begin{center}
  {\color{white}\textbf{\LARGE AP Precalculus}}\\[4pt]
  {\color{skymid}\large\textbf{Unit 2: Exponential and Logarithmic Functions}}\\[6pt]
  {\color{white}\textbf{\large Lesson 2.15 \quad Semi-log Plots}}
\end{center}

\vspace{0.22in}
\namedateperiod
\vspace{0.18in}

\begin{learningtargetbox}
\small
\begin{enumerate}[leftmargin=1.5em, itemsep=5pt]
  \item I can explain what a semi-log plot is and why exponential data appears linear on one.
        \hfill\textit{LO 2.15.A, EK 2.15.A.1}
  \item I can determine whether an exponential model is appropriate by examining whether
        $\log_{10}(y)$ vs.\ $x$ forms a linear pattern. \hfill\textit{LO 2.15.A}
  \item I can construct the linearization of exponential data using
        $\log_n(y) = (\log_n b)\cdot x + \log_n a$. \hfill\textit{LO 2.15.B, EK 2.15.B.2}
  \item I can recover the exponential model $y = ab^x$ from the slope and intercept of the
        linearized data. \hfill\textit{LO 2.15.B, EK 2.15.B.1}
\end{enumerate}
\end{learningtargetbox}

\vspace{0.14in}

\begin{tocbox}
\small
\renewcommand{\arraystretch}{1.5}
\begin{tabularx}{\linewidth}{@{} c l X r @{}}
\rowcolor{sky}
\textbf{\#} & \textbf{Component} & \textbf{Description} & \textbf{Score} \\
\midrule
1 & Warm-Up        & Spiral review: exponential values, $\log_{10}$, and constant rate of change & \blank{1.2cm} \\
2 & Guided Notes   & Semi-log plots; linearizing $y = ab^x$; recovering the exponential model & \blank{1.2cm} \\
3 & Group Activity & Drug-concentration data; building and verifying exponential models & \blank{1.2cm} \\
4 & Exit Ticket    & Compute $\log_{10}(y)$; find slope and intercept; write exponential model & \blank{1.2cm} \\
5 & Homework       & Independent practice with semi-log linearization and model recovery & \blank{1.2cm} \\
\midrule
  & \multicolumn{2}{r}{\textbf{Total}} & \blank{1.2cm} \\
\end{tabularx}
\end{tocbox}

\end{document}
